\section{State where the stream begins and ends}
This codec receives its byte extent from the caller. Continuation bits locate value
boundaries inside that extent, but they do not identify the expected number of readings.
If a stream is cut immediately after a complete value, the prefix is another valid
shorter stream. Detecting a missing complete suffix requires framing information such
as an expected count or an external record length.

A call also resets the predecessor to zero. Encoding the second half of a sequence in
a fresh call produces a new first delta relative to zero. Concatenating independently
encoded calls is therefore not generally equivalent to encoding their combined input.
A larger file format can define separate frames with reset points, or a future interface
can pass the predecessor explicitly. Do not silently carry state between calls to make
one example resume.

Compression provides no integrity check. A changed payload can remain canonical and
reconstruct a different in-range reading. A checksum can detect some accidental changes
when specified by an enclosing format; authenticity needs a separate design. The next
project develops a file checksum, while this chapter keeps its own representation and
error contract explicit.

\Needspace{10\baselineskip}
\begin{MARSExercise}{Define a resumable extension}
  The encoder stopped because the last complete value did not fit. Design an additional
  argument or result for resuming from the unconsumed reading in a fresh call. State
  which predecessor must be transferred and whether the decoder needs matching frame
  metadata. Keep this proposed extension separate from the current interface.
\end{MARSExercise}
